% GENERATED FILE. Edit canonical lessons, then run npm run generate:books.
% canonical-source-hash: fnv1a64:474ef67fa6d86d84
% canonical-lessons: GE-C167-bequem, GE-C167-unfreundlich, GE-C167-hoeflich, GE-C167-spannend, GE-C167-moeglich, GE-R167-first-pass-words-from-chapters-162-164, GE-R167-second-pass-words-from-chapters-165-167

\chapter{Comfortable, Unfriendly, Polite, Exciting, Possible}
\label{ch:ge-comfortable-unfriendly-polite-exciting-possible}
\hlchaptermodality{167}

\begin{chapteropening}
\textbf{By the end of this chapter:} \emph{I can say comfortable, unfriendly, polite, exciting and possible.}

Complete the last lesson of chapter 167: i can say comfortable, unfriendly, polite, exciting and possible.
\end{chapteropening}

\section[bequem]{bequem — comfortable}

\label{lesson:GE-C167-bequem}



\begin{quote}
Before the new one: say the German for a possibility, then the German for the afternoon.
\end{quote}

\subsection*{You'll want to know: bequem}
\textbf{bequem} — \textquotedblleft{}comfortable\textquotedblright{}.

\begin{grammarlens}[title={What you've built}]
One more word for this chapter.
\end{grammarlens}

\subsection*{Your turn}
\begin{itemize}
\raggedright
  \item \emph{Say it:} \emph{bequem}
  \item \emph{Say it:} \emph{bequem}, once more
  \item \emph{Say it:} say \emph{bequem}
  \item \emph{Recall:} say the German for a possibility, then the German for the afternoon, then say \emph{bequem} again
\end{itemize}

\begin{tcolorbox}[breakable,colback=teal!4,colframe=teal!35!black,title={Before you move on}]
What does bequem mean? (\textquotedblleft{}Comfortable\textquotedblright{}.) Say it once more.
\end{tcolorbox}
\section[unfreundlich]{unfreundlich — unfriendly}

\label{lesson:GE-C167-unfreundlich}



\begin{quote}
Before the new one: say the German for the afternoon, then the German for comfortable.
\end{quote}

\subsection*{You'll want to know: unfreundlich}
\textbf{unfreundlich} — \textquotedblleft{}unfriendly\textquotedblright{}.

\begin{grammarlens}[title={What you've built}]
One more word for this chapter.
\end{grammarlens}

\subsection*{Your turn}
\begin{itemize}
\raggedright
  \item \emph{Say it:} \emph{unfreundlich}
  \item \emph{Say it:} \emph{unfreundlich}, once more
  \item \emph{Say it:} say \emph{unfreundlich}
  \item \emph{Recall:} say the German for the afternoon, then the German for comfortable, then say \emph{unfreundlich} again
\end{itemize}

\begin{tcolorbox}[breakable,colback=teal!4,colframe=teal!35!black,title={Before you move on}]
What does unfreundlich mean? (\textquotedblleft{}Unfriendly\textquotedblright{}.) Say it once more.
\end{tcolorbox}
\section[höflich]{höflich — polite}

\label{lesson:GE-C167-hoeflich}



\begin{quote}
Before the new one: say the German for comfortable, then the German for unfriendly.
\end{quote}

\subsection*{You'll want to know: höflich}
\textbf{höflich} — \textquotedblleft{}polite\textquotedblright{}.

\begin{grammarlens}[title={What you've built}]
One more word for this chapter.
\end{grammarlens}

\subsection*{Your turn}
\begin{itemize}
\raggedright
  \item \emph{Say it:} \emph{höflich}
  \item \emph{Say it:} \emph{höflich}, once more
  \item \emph{Say it:} say \emph{höflich}
  \item \emph{Recall:} say the German for comfortable, then the German for unfriendly, then say \emph{höflich} again
\end{itemize}

\begin{tcolorbox}[breakable,colback=teal!4,colframe=teal!35!black,title={Before you move on}]
What does höflich mean? (\textquotedblleft{}Polite\textquotedblright{}.) Say it once more.
\end{tcolorbox}
\section[spannend]{spannend — exciting}

\label{lesson:GE-C167-spannend}



\begin{quote}
Before the new one: say the German for unfriendly, then the German for polite.
\end{quote}

\subsection*{You'll want to know: spannend}
\textbf{spannend} — \textquotedblleft{}exciting\textquotedblright{}.

\begin{grammarlens}[title={What you've built}]
One more word for this chapter.
\end{grammarlens}

\subsection*{Your turn}
\begin{itemize}
\raggedright
  \item \emph{Say it:} \emph{spannend}
  \item \emph{Say it:} \emph{spannend}, once more
  \item \emph{Say it:} say \emph{spannend}
  \item \emph{Recall:} say the German for unfriendly, then the German for polite, then say \emph{spannend} again
\end{itemize}

\begin{tcolorbox}[breakable,colback=teal!4,colframe=teal!35!black,title={Before you move on}]
What does spannend mean? (\textquotedblleft{}Exciting\textquotedblright{}.) Say it once more.
\end{tcolorbox}
\section[möglich]{möglich — possible}

\label{lesson:GE-C167-moeglich}



\begin{quote}
Before the new one: say the German for polite, then the German for exciting.
\end{quote}

\subsection*{You'll want to know: möglich}
\textbf{möglich} — \textquotedblleft{}possible\textquotedblright{}.

\begin{grammarlens}[title={What you've built}]
One more word for this chapter.
\end{grammarlens}

\subsection*{Your turn}
\begin{itemize}
\raggedright
  \item \emph{Say it:} \emph{möglich}
  \item \emph{Say it:} \emph{möglich}, once more
  \item \emph{Say it:} say \emph{möglich}
  \item \emph{Recall:} say the German for polite, then the German for exciting, then say \emph{möglich} again
\end{itemize}

\begin{tcolorbox}[breakable,colback=teal!4,colframe=teal!35!black,title={Before you move on}]
What does möglich mean? (\textquotedblleft{}Possible\textquotedblright{}.) Say it once more.
\end{tcolorbox}
\section[(dialogue)]{Words from chapters 162-164 — a first pass}

\label{lesson:GE-R167-first-pass-words-from-chapters-162-164}



\begin{quote}
Say the words for exciting and possible.
\end{quote}

\subsection*{Your turn}
Say these aloud:

\begin{itemize}
\raggedright
  \item a sofa, a radio, the weather, the wind, a thunderstorm
  \item the temperature, a birthday, a wedding, a party, an invitation
  \item a neighbour, a cousin, a granddaughter, a baby, a human being
\end{itemize}

\begin{tcolorbox}[breakable,colback=teal!4,colframe=teal!35!black,title={Before you move on}]
A sofa? (\textbf{das Sofa}.) A radio? (\textbf{das Radio}.) The weather? (\textbf{das Wetter}.) The wind? (\textbf{der Wind}.) A thunderstorm? (\textbf{das Gewitter}.) The temperature? (\textbf{die Temperatur}.) A birthday? (\textbf{der Geburtstag}.) A wedding? (\textbf{die Hochzeit}.) A party? (\textbf{die Party}.) An invitation? (\textbf{die Einladung}.) A neighbour? (\textbf{die Nachbarin}.) A cousin? (\textbf{die Cousine}.) A granddaughter? (\textbf{die Enkelin}.) A baby? (\textbf{das Baby}.) A human being? (\textbf{der Mensch}.)
\end{tcolorbox}
\section[(dialogue)]{Words from chapters 165-167 — a second pass}

\label{lesson:GE-R167-second-pass-words-from-chapters-165-167}



\begin{quote}
Say the words for people and luck.
\end{quote}

\subsection*{Your turn}
Say these aloud:

\begin{itemize}
\raggedright
  \item people, luck, fun, desire, an opinion
  \item a reason, a solution, experience, a possibility, the afternoon
  \item comfortable, unfriendly, polite, exciting, possible
\end{itemize}

\begin{tcolorbox}[breakable,colback=teal!4,colframe=teal!35!black,title={Before you move on}]
People? (\textbf{die Leute}.) Luck? (\textbf{das Glück}.) Fun? (\textbf{der Spaß}.) Desire? (\textbf{die Lust}.) An opinion? (\textbf{die Meinung}.) A reason? (\textbf{der Grund}.) A solution? (\textbf{die Lösung}.) Experience? (\textbf{die Erfahrung}.) A possibility? (\textbf{die Möglichkeit}.) The afternoon? (\textbf{der Nachmittag}.) Comfortable? (\textbf{bequem}.) Unfriendly? (\textbf{unfreundlich}.) Polite? (\textbf{höflich}.) Exciting? (\textbf{spannend}.) Possible? (\textbf{möglich}.)
\end{tcolorbox}
